\documentclass[10pt,a4paper]{amsart}
\usepackage[margin=19mm]{geometry}
\usepackage{amsmath,amssymb,mathtools}
\usepackage[scr=rsfs,cal=euler]{mathalfa}
\usepackage{xfrac}
\usepackage[unicode,hidelinks,bookmarks=false]{hyperref}
\newcommand{\ie}{i.e.\ }
\newcommand{\cf}{cf.\ }
\newcommand{\resp}{respectively\ }
\newcommand{\Subsection}{\subsection}
\providecommand{\nobreakdash}{\nobreak\hbox{-}}
\setlength{\emergencystretch}{2em}
\multlinegap0pt
\usepackage{amssymb}
\usepackage{mathtools}
\usepackage{bm}
\usepackage[shortlabels]{enumitem}
\usepackage{xcolor}
\newtheorem{lemma}{Lemma}
\title{\bf \Large}
\author{Yongjiang Wu, Lihua Feng}
\date{Source: arXiv:2607.26871v1 (CC BY 4.0)}
\begin{document}
\maketitle

\noindent\emph{Source and licence.} \bf \Large. Version arXiv:2607.26871v1 (CC BY 4.0), distributed by arXiv under the Creative Commons Attribution 4.0 International licence, http://creativecommons.org/licenses/by/4.0/, as recorded in the arXiv licence metadata for this exact version and shown on the abstract page https://arxiv.org/abs/2607.26871v1. Full licence notice: \texttt{licenses/arxiv-2607.26871-CC-BY-4.0.txt}.\par\smallskip

\noindent\textit{Editorial scope.} Counted unit: the article's lemma lem:bellman (source0.tex lines 1111--1137). The article's complete original proof of it is delivered with it (source0.tex lines 1139--1968, one \texttt{proof} environment of 830 lines). No mathematics is written, repaired or re-derived: the statement and the proof are the article's own lines, in the article's own notation, and no retained line is substituted or re-flowed.  No further source-local context is needed: every cross-reference of every retained line resolves inside the two delivered spans.  Nothing was omitted from any retained span, so the package carries no omission record.  No retained line contains a citation, so the wrapper carries no bibliography and no bibliography entry.  No retained line contains a figure environment, an image-inclusion command or drawing code, so no image is delivered and the package declares no asset dependency.  Editorial formatting only, in addition: an AMS \texttt{amsart} preamble that restates the article's own declarations and macros used by the retained lines, namely the environments \texttt{lemma} on separate counters, together with the standard packages those lines need.  The retained environments print in this document as Lemma 1 in order of appearance, whereas the article prints the same items with the numbering of its own full document.  The complete native source of the article as distributed in the arXiv e-print for this exact version is bundled unchanged as \texttt{sources/206300/source0.tex}.
\begin{lemma}\label{lem:bellman}
Let $K>1$ and suppose that
\begin{equation*}
 \frac{1}{2(K+1)}
 \leq d_1
 \leq
 \frac{1}{(1+\sqrt K)^2}.
\end{equation*}
Let $0\leq x\leq y\leq1$ and $0\leq v\leq u\leq1$ satisfy
\begin{equation*}
 xu\leq d_1,
 \qquad
 yv\leq d_1,
 \qquad
 u\leq h_K(x),
 \qquad
 v\leq h_K(y).
\end{equation*}
Then, for $0<p,q\leq\frac{1}{2}$,
\begin{equation}
 \bigl((1-p)x+py\bigr)
 \bigl(qu+(1-q)v\bigr)
 \leq
 \max\{d_1,pq\}.
 \label{eq:bellman}
\end{equation}
\end{lemma}

\begin{proof}
Define
\begin{equation*}
 F(t)=t h_K(t)
     =\frac{t(1-t)}{1+(K-1)t},
 \qquad 0\leq t\leq1.
\end{equation*}
Differentiating gives
\begin{equation*}
 F'(t)
 =
 \frac{1-2t-(K-1)t^2}
      {(1+(K-1)t)^2}.
\end{equation*}
The numerator has a unique zero in $(0,1)$, namely
\begin{equation*}
 t_0=\frac{1}{1+\sqrt K}.
\end{equation*}
Consequently, $F$ is increasing on $[0,t_0]$ and decreasing on
$[t_0,1]$. Moreover,
\begin{equation*}
 F(t_0)=\frac{1}{(1+\sqrt K)^2},
 \qquad
 F\left(\frac12\right)=\frac{1}{2(K+1)}.
\end{equation*}
The assumed range of $d_1$ therefore implies that the equation
$
 t h_K(t)=d_1
$
has two roots, counted with multiplicity,
\begin{equation*}
 0<a\leq t_0\leq b\leq\frac12.
\end{equation*}

The equation $t h_K(t)=d_1$ can be rewritten as
$
 t^2-\bigl(1-(K-1)d_1\bigr)t+d_1=0.
$
Hence,
\begin{equation*}
 ab=d_1,
 \qquad
 a+b=1-(K-1)d_1.
\end{equation*}
Set
$
 c:=1-a-b=(K-1)d_1>0.
$
Then
\begin{equation*}
 K=\frac{(1-a)(1-b)}{ab},\qquad
 h_K(t)
 =
 \frac{d_1(1-t)}{d_1+ct}.
\end{equation*}

For $t>0$, define
\begin{equation*}
 \phi(t)=\min\left\{h_K(t),\frac{d_1}{t}\right\},
\end{equation*}
and set $\phi(0)=1$. Since
\begin{equation*}
 h_K(t)-\frac{d_1}{t}
 =
 -\frac{d_1(t-a)(t-b)}
        {t(d_1+ct)},
\end{equation*}
we have
\begin{equation}
 \phi(t)=
 \begin{cases}
  h_K(t),&0\leq t\leq a,\\[2pt]
  d_1/t,&a\leq t\leq b,\\[2pt]
  h_K(t),&b\leq t\leq1.
 \end{cases}
 \label{eq:phi-pieces}
\end{equation}
A direct calculation gives
\begin{equation*}
 h_K(h_K(t))=t.
\end{equation*}
Furthermore,
\begin{equation*}
 h_K(a)=b,
 \qquad
 h_K(b)=a.
\end{equation*}
Since $h_K$ is decreasing and
\begin{equation*}
 h_K(0)=1,
 \qquad
 h_K(a)=b,
 \qquad
 h_K(b)=a,
 \qquad
 h_K(1)=0,
\end{equation*}
we have
\begin{equation*}
 h_K([0,a])=[b,1],
 \qquad
 h_K([b,1])=[0,a].
\end{equation*}
Moreover, the map $t\mapsto d_1/t$ maps $[a,b]$ onto itself. Both maps are decreasing
and involutive on the corresponding intervals. It follows that
$\phi$ is decreasing and satisfies
\begin{equation*}
 \phi(\phi(t))=t
 \qquad\text{for every }t\in[0,1].
\end{equation*}

The inequalities $xu\leq d_1$ and $u\leq h_K(x)$ imply
$
 u\leq\phi(x).
$
Indeed, this is immediate when $x=0$, while for $x>0$ we have both
$u\leq h_K(x)$ and $u\leq d_1/x$. Similarly,
$
 v\leq\phi(y).
$
Since $\phi$ is decreasing and involutive, the latter inequality gives
\begin{equation*}
 \phi(v)\geq\phi(\phi(y))=y.
\end{equation*}
Moreover, $x\leq y$ implies
\begin{equation*}
 v\leq\phi(y)\leq\phi(x).
\end{equation*}
The left-hand side of \eqref{eq:bellman} is nondecreasing in $u$ and
$y$. We may therefore replace $u$ by $\phi(x)$ and $y$ by $\phi(v)$.
It remains to prove
\begin{equation*}
 G_{p,q}(x,v)\leq\max\{d_1,pq\},
\end{equation*}
where
\begin{equation*}
 G_{p,q}(x,v)
 =
 \bigl((1-p)x+p\phi(v)\bigr)
 \bigl((1-q)v+q\phi(x)\bigr),\qquad v\leq\phi(x).
\end{equation*}
Since $\phi$ is a decreasing involution, the last inequality is
equivalent to $x\leq\phi(v)$.


We first reduce the proof to the case $pq=d_1$. Observe that $d_1<1/4$. The inequalities
$\phi(v)\geq x$ and $\phi(x)\geq v$ give
\begin{align*}
 \frac{\partial G_{p,q}}{\partial p}
 &=
 \bigl(\phi(v)-x\bigr)
 \bigl((1-q)v+q\phi(x)\bigr)
 \geq0,\\
 \frac{\partial G_{p,q}}{\partial q}
 &=
 \bigl(\phi(x)-v\bigr)
 \bigl((1-p)x+p\phi(v)\bigr)
 \geq0.
\end{align*}
Thus, $G_{p,q}$ is nondecreasing in each of $p$ and $q$. On the other
hand,
\begin{equation}
 \frac{G_{p,q}(x,v)}{pq}
 =
 \bigl(\phi(v)+(p^{-1}-1)x\bigr)
 \bigl(\phi(x)+(q^{-1}-1)v\bigr).\label{eq:normalized-G}
\end{equation}
For fixed $x$ and $v$, the right-hand side of
\eqref{eq:normalized-G} is nonincreasing separately in $p$ and $q$.




Suppose first that $pq\leq d_1$. For $0\leq t\leq1$, define
\begin{equation*}
 p(t)=p+t\left(\frac12-p\right),
 \qquad
 q(t)=q+t\left(\frac12-q\right).
\end{equation*}
Then $p(t)$ and $q(t)$ are nondecreasing functions of $t$, and
\begin{equation*}
 p(0)q(0)=pq\leq d_1,
 \qquad
 p(1)q(1)=\frac14>d_1.
\end{equation*}
Since $p(t)q(t)$ is continuous, there exists $t_0\in[0,1)$ such that
$
 p(t_0)q(t_0)=d_1.
$
Set
$
 p'=p(t_0)
$
and 
$
 q'=q(t_0).
$
Then
\begin{equation*}
 p'\geq p,
 \qquad
 q'\geq q,
 \qquad
 p'q'=d_1.
\end{equation*}
Since $G_{p,q}$ is nondecreasing separately in $p$ and $q$, we
have
$
 G_{p,q}(x,v)
 \leq
 G_{p',q'}(x,v).
$
Thus, it suffices to prove
\begin{equation*}
 G_{p',q'}(x,v)\leq d_1
\end{equation*}
Therefore, the case $pq\leq d_1$ is reduced to the case $pq=d_1$.


Suppose next that $pq\geq d_1$. Keep $p$ fixed and define
$
 q'=\frac{d_1}{p}.
$
Then $q'>0$ and
$
 q'\leq q\leq\frac12.
$
By construction,
$
 pq'=d_1.
$
For fixed $p,x$ and $v$, equation \eqref{eq:normalized-G} shows that
\begin{equation*}
 \frac{G_{p,s}(x,v)}{ps}
\end{equation*}
is nonincreasing as a function of $s$. Since $q'\leq q$, it follows
that
\begin{equation*}
 \frac{G_{p,q}(x,v)}{pq}
 \leq
 \frac{G_{p,q'}(x,v)}{pq'}.
\end{equation*}
Once the estimate has been proved for pairs whose product is $d_1$, we
may apply it to $(p,q')$ to obtain
$
 G_{p,q'}(x,v)\leq d_1=pq'.
$
Thus,
\begin{equation*}
 \frac{G_{p,q}(x,v)}{pq}
 \leq
 \frac{G_{p,q'}(x,v)}{pq'}
 \leq1,
\end{equation*}
and hence
\begin{equation*}
 G_{p,q}(x,v)
 \leq pq
 =
 \max\{d_1,pq\}.
\end{equation*}
Therefore, the case $pq\geq d_1$ is also reduced to the case $pq=d_1$.


















\textbf{It remains to prove
$
 G_{p,q}(x,v)\leq d_1
$
under the assumption
$
 pq=d_1.
$}
In this case, we have
$
 \frac{d_1}{p}=q\leq\frac12,
$
which is equivalent to
$
 p\geq2d_1.
$
Together with $p\leq1/2$, this gives
\begin{equation*}
 2d_1\leq p\leq\frac12.
\end{equation*}
By symmetry, $p\leq1/2$ and $pq=d_1$ also give
\begin{equation*}
 2d_1\leq q\leq\frac12.
\end{equation*}
Consequently, it suffices to prove $
 G_{p,q}(x,v)\leq d_1
$ under
\begin{equation*}
 pq=d_1,
 \qquad
 2d_1\leq p,q\leq\frac12.
\end{equation*}



Recall that $0<a\leq b\leq\frac12$.
Let
$
 I_0=[0,a],
 I_1=[a,b]
$
and
$
 I_2=[b,1].
$
By \eqref{eq:phi-pieces} and the involution property of $\phi$, we have
\begin{equation*}
 \phi(I_0)=I_2,
 \qquad
 \phi(I_1)=I_1,
 \qquad
 \phi(I_2)=I_0.
\end{equation*}
We now determine which pairs of intervals can contain $(x,v)$ under
the constraint $v\leq\phi(x)$.

If $x\in I_0$, then $\phi(x)\in I_2$, so $v$ may lie in $I_0$,
$I_1$, or $I_2$. If $x\in I_1$, then $\phi(x)\in I_1$ and hence
$v\leq b$, so $v$ can lie only in $I_0$ or $I_1$. Finally, if
$x\in I_2$, then $\phi(x)\in I_0$ and hence $v\leq a$, so
$v$ must lie in $I_0$.
Consequently, every pair $(x,v)$ satisfying $v\leq\phi(x)$ belongs to
one of the following six regions:
\begin{equation*}
 I_0\times I_0,\quad
 I_0\times I_1,\quad
 I_0\times I_2,\quad
 I_1\times I_0,\quad
 I_1\times I_1,\quad
 I_2\times I_0.
\end{equation*}
Within each of these regions, the constraint $v\leq\phi(x)$ remains in
force.


\textbf{Suppose first that $(x,v)\in I_1\times I_1$.} Then
$
 \phi(x)=\frac{d_1}{x}
$
and 
$
 \phi(v)=\frac{d_1}{v}.
$
Put $z=xv$. Since $x,v\geq a>0$ and $v\leq \phi(x)=d_1/x$, we have
\begin{equation*}
 a^2\leq z\leq d_1.
\end{equation*}
Expanding $G_{p,q}$ and using $pq=d_1$ gives
\begin{equation*}
 G_{p,q}(x,v)
 =
 (1-p)(1-q)z
 +
 \bigl((1-p)q+p(1-q)\bigr)d_1
 +
 \frac{d_1^3}{z}.
\end{equation*}
Denote the right-hand side by
$
 H(z).
$
Since $z>0$, we have
\begin{equation*}
 H''(z)=\frac{2d_1^3}{z^3}>0.
\end{equation*}
Thus, $H$ is convex on $[a^2,d_1]$. For any
$z\in[a^2,d_1]$, write
$
 z=(1-\theta)a^2+\theta d_1
$
for some $\theta\in[0,1]$. Convexity gives
\begin{equation*}
 H(z)
 \leq
 (1-\theta)H(a^2)+\theta H(d_1)
 \leq
 \max\{H(a^2),H(d_1)\}.
\end{equation*}
Thus, it suffices to verify the desired inequality at
$z=a^2$ and $z=d_1$. 
At $z=d_1$, we have $xv=d_1$. Since $x,v\in I_1$,
\begin{equation*}
 \phi(x)=\frac{d_1}{x}=v,
 \qquad
 \phi(v)=\frac{d_1}{v}=x.
\end{equation*}
Therefore,
\begin{align*}
 G_{p,q}(x,v)
 =
 \bigl((1-p)x+p\phi(v)\bigr)
 \bigl((1-q)v+q\phi(x)\bigr)=xv=d_1.
\end{align*} At $z=a^2$, the inequalities $x,v\geq a$ and the identity
$xv=a^2$ imply that $x=v=a$. Since $d_1=ab$, we have
\begin{equation*}
 \phi(a)=\frac{d_1}{a}=b.
\end{equation*}
Therefore,
\begin{align*}
 G_{p,q}(a,a)
 =
 \bigl((1-p)a+p\phi(a)\bigr)
 \bigl((1-q)a+q\phi(a)\bigr)=
 \bigl(a+p(b-a)\bigr)
 \bigl(a+q(b-a)\bigr).
\end{align*}
Since
$
 (1-2p)(1-2q)\geq0,
$
we have
\begin{equation*}
 p+q\leq\frac12+2pq=\frac12+2ab.
\end{equation*}
Consequently,
\begin{align*}
 d_1-
 \bigl(a+p(b-a)\bigr)
 \bigl(a+q(b-a)\bigr)&=
 a(b-a)\bigl(1-p-q-b(b-a)\bigr)\\
 &\geq
 a(b-a)\left(\frac12-b(a+b)\right)\geq0.
\end{align*}
where the last inequality follows from $b\leq1/2$ and $a+b\leq1$.
Therefore, $G_{p,q}(x,v)\leq d_1$ on $I_1\times I_1$.

\textbf{We next consider $(x,v)\in I_0\times I_1$.} Here,
$
 \phi(x)=h_K(x)
$
and
$
 \phi(v)=\frac{d_1}{v}.
$
Recall that $G_{p,q}(x,v)
 =
 \bigl((1-p)x+p\phi(v)\bigr)
 \bigl((1-q)v+q\phi(x)\bigr).$
Since $x\leq a\leq1/2<1$, we have $h_K(x)>0$.
For fixed $x$, differentiation with respect to $v$ gives
\begin{equation*}
 \frac{\partial^2G_{p,q}}{\partial v^2}
 =
 \frac{2d_1^2h_K(x)}{v^3}>0.
\end{equation*}
Thus, for each fixed $x\in I_0$, the function
$v\mapsto G_{p,q}(x,v)$ is convex on $[a,b]$. As in the preceding
convexity argument, this implies
\begin{equation*}
 G_{p,q}(x,v)
 \leq
 \max\bigl\{G_{p,q}(x,a),G_{p,q}(x,b)\bigr\}
 \qquad\text{for every }v\in[a,b].
\end{equation*}
Therefore, it suffices to verify the desired inequality at
$v=a$ and $v=b$.

Put
$
 D_x=d_1+cx.
$
Substituting $q=d_1/p$ and
$h_K(x)=d_1(1-x)/D_x$, and then collecting powers of $x$, gives
\begin{align}
 d_1-G_{p,q}(x,a)
 &=
 \frac{a(b-x)Q_a(x)}{pD_x},
 \label{eq:Qa-factor}\\
 d_1-G_{p,q}(x,b)
 &=
 \frac{b(a-x)Q_b(x)}{pD_x},
 \label{eq:Qb-factor}
\end{align}
where $Q_a$ and $Q_b$ are the following affine functions:
\begin{align*}
 Q_a(x)
 &=
 \left(1-\frac{x}{a}\right)
 abp(1-b-p+ab)
 +
 \frac{x}{a}\,a(a-1)R_a(p),\\
 Q_b(x)
 &=
 \left(1-\frac{x}{a}\right)
 abp(1-a-p+ab)
 +
 \frac{x}{a}\,aR_b(p),
\end{align*}
with
\begin{align*}
 R_a(p)
 &=
 p^2+p(b^2-ab-1)+ab,\\
 R_b(p)
 &=
 (1-a)p(1-p)-ab(1-b).
\end{align*}

Since $pq=d_1=ab$ and $q\leq1/2$, we have
$p=ab/q\geq2ab$.
Since $0<a,b,p\leq1/2$, we have
\begin{equation*}
 Q_a(0)=abp(1-b-p+ab)\geq0,\qquad
 Q_b(0)=abp(1-a-p+ab)\geq0.
\end{equation*}
The polynomial $R_a(p)$ is convex, and
\begin{align*}
 R_a(2ab)
 =
 ab\bigl(2b(a+b)-1\bigr)\leq0,\qquad
 R_a\left(\frac12\right)
 =
 \frac{2b(a+b)-1}{4}\leq0.
\end{align*}
As in the preceding
convexity argument, $R_a(p)\leq0$ for $2ab\leq p\leq1/2$. Since
$a(a-1)<0$, it follows that
\begin{equation*}
 Q_a(a)=a(a-1)R_a(p)\geq0.
\end{equation*}
As $Q_a$ is affine, $Q_a(x)\geq0$ for every $x\in[0,a]$.
Similarly, $R_b(p)$ is concave, and
\begin{align*}
 R_b(2ab)
 =
 ab(2a-1)(2ab-b-1)\geq0,\qquad
 R_b\left(\frac12\right)
 =
 \frac{1-2a+a(1-2b)^2}{4}\geq0.
\end{align*}
Hence, $R_b(p)\geq0$ for $2ab\leq p\leq1/2$, and therefore
\begin{equation*}
 Q_b(a)=aR_b(p)\geq0.
\end{equation*}
Because $Q_b$ is affine, $Q_b(x)\geq0$ on $[0,a]$.
Equations \eqref{eq:Qa-factor} and \eqref{eq:Qb-factor} now show that
\begin{equation*}
 G_{p,q}(x,v)\leq d_1
 \qquad\text{on }I_0\times I_1.
\end{equation*}
Since $pq=d_1$ and $p\leq1/2$, we also have
$q=d_1/p\geq2d_1$. Hence, $2d_1\leq q\leq1/2$.
In  view of  $G_{p,q}(x,v)=G_{q,p}(v,x)$
and $2d_1\leq q\leq1/2$, \textbf{the same argument proves the result on
$I_1\times I_0$}. Indeed, the constraint $v\leq\phi(x)$ is equivalent to
$x\leq\phi(v)$. Thus,
after interchanging $(x,p)$ with $(v,q)$, a feasible point in
$I_1\times I_0$ becomes a feasible point in $I_0\times I_1$. 


\textbf{It remains to treat the three regions on which both occurrences of
$\phi$ are equal to $h_K$.} Write
\begin{equation*}
 A=(1-p)x+ph_K(v),
 \qquad
 B=(1-q)v+qh_K(x),
 \qquad
 G=AB,\qquad  D_t=d_1+ct.
\end{equation*}
Observe that $G= G_{p,q}(x,v)$, and
\begin{equation*}
 h_K'(t)
 =
 -\frac{d_1(d_1+c)}{D_t^2},
 \qquad
 h_K''(t)
 =
 \frac{2cd_1(d_1+c)}{D_t^3}.
\end{equation*}
Consequently,
\begin{align*}
 G_x=(1-p)B+qh_K'(x)A,\qquad G_v=ph_K'(v)B+(1-q)A,
\end{align*}
and
\begin{align*}
 G_{xx}
 &=
 \frac{2qd_1(d_1+c)}{D_x^3}
 \bigl(cp h_K(v)-d_1(1-p)\bigr),\\
 G_{vv}
 &=
 \frac{2pd_1(d_1+c)}{D_v^3}
 \bigl(cq h_K(x)-d_1(1-q)\bigr).
\end{align*}
Here and below, subscripts on $G$ denote partial derivatives with
respect to the indicated variables.
Suppose that $G$ attains a local maximum at an
interior point $(x,v)$ of one of these regions. Then
\begin{equation*}
 G_x=G_v=0,\qquad
 G_{xx}\leq0,
 \qquad
 G_{vv}\leq0,
 \qquad
 G_{xx}G_{vv}-G_{xv}^2\geq0.
\end{equation*}
At an interior point, $x,v<1$, and hence
$h_K(x),h_K(v)>0$. Since $p,q>0$, we have
\begin{equation*}
 A\geq ph_K(v)>0,
 \qquad
 B\geq qh_K(x)>0.
\end{equation*}
The equations $G_x=G_v=0$ imply
\begin{equation*}
 pqh_K'(x)h_K'(v)
 =
 (1-p)(1-q)=:L.
\end{equation*}
Moreover,
\begin{equation*}
 G_{xv}
 =
 (1-p)(1-q)+pqh_K'(x)h_K'(v)
 =
 2L>0.
\end{equation*}
Define
\begin{equation*}
 \Delta_p=d_1(1-p)-cp h_K(v),
 \qquad
 \Delta_q=d_1(1-q)-cq h_K(x).
\end{equation*}
Since $ G_{xx}\leq0$ and $
 G_{vv}\leq0$, we have $\Delta_p,\Delta_q\geq0$. If $\Delta_p=0$, then $G_{xx}=0$, while if $\Delta_q=0$, then
$G_{vv}=0$. In either case, since $G_{xv}>0$, we have
\begin{equation*}
 G_{xx}G_{vv}-G_{xv}^2
 =
 -G_{xv}^2<0,
\end{equation*}
a contradiction. Therefore, we have
$
 \Delta_p, \Delta_q>0.
$
Since $pq=d_1$, $G_{xv}=2L$ and
\begin{equation*}
 L=pqh_K'(x)h_K'(v)
  =\frac{d_1^3(d_1+c)^2}{D_x^2D_v^2},
\end{equation*}
the formulas for $G_{xx}$ and $G_{vv}$ give
\begin{align*}
 \frac{G_{xx}G_{vv}}{G_{xv}^2}
 =
 \frac{\Delta_p\Delta_q}{LD_xD_v}<
 \frac{d_1^2(1-p)(1-q)}{LD_xD_v}=
 \frac{d_1^2}{D_xD_v}
 \leq1.
\end{align*}
Thus,
\begin{equation*}
 G_{xx}G_{vv}-G_{xv}^2<0,
\end{equation*}
again a contradiction. Therefore, none of
these three regions has an interior local maximum.

It remains to check their boundaries. We first consider the boundary $v=0$. Since $h_K(0)=1$, we have
\begin{equation*}
 G(x,0)
 =
 \bigl(p+(1-p)x\bigr)qh_K(x).
\end{equation*}
Using
\begin{equation*}
 q=\frac{d_1}{p},
 \qquad
 h_K(x)=\frac{d_1(1-x)}{D_x},
\end{equation*}
we obtain
\begin{equation*}
 G(x,0)
 =
 \frac{d_1^2(1-x)\bigl(p+(1-p)x\bigr)}
 {pD_x}.
\end{equation*}
Consequently,
\begin{align*}
 d_1-G(x,0)
 &=
 \frac{d_1}{pD_x}
 \left[
 pD_x-d_1(1-x)\bigl(p+(1-p)x\bigr)
 \right]\\
&=\frac{d_1x}{pD_x}
 \bigl(pc-d_1(1-2p)+d_1(1-p)x\bigr).
\end{align*}
Recall that $d_1=pq$ and
$
 c=1-a-b.
$
Since $q\leq1/2$, we have
\begin{align*}
 q(1-2p)
 =q-2pq=q-2d_1\leq\frac12-2d_1.
\end{align*}
Moreover,
\begin{equation*}
 c-\left(\frac12-2d_1\right)
 =
 2\left(\frac12-a\right)
  \left(\frac12-b\right)
 \geq0.
\end{equation*}
Therefore,
$
 q(1-2p)\leq c.
$
Multiplying by $p>0$ gives
$
 d_1(1-2p)=pq(1-2p)\leq pc.
$
Thus,
\begin{equation*}
 pc-d_1(1-2p)+d_1(1-p)x\geq0.
\end{equation*}
It follows that
\begin{equation*}
 d_1-G(x,0)\geq0.
\end{equation*}


On $x=0$, direct simplification gives
\begin{equation*}
 d_1-G(0,v)
 =
 \frac{d_1v
 \bigl(c+2d_1-p+v(p-d_1)\bigr)}
 {d_1+cv}.
\end{equation*}
Observe that
$
 p-d_1=p(1-q)>0
$
and
\begin{align*}
 c+2d_1-p
 \geq
 c+2d_1-\frac12=
 2\left(\frac12-a\right)
  \left(\frac12-b\right)
 \geq0.
\end{align*}
Therefore, $G(0,v)\leq d_1$. 
The edge $v=a$ was treated in \eqref{eq:Qa-factor}, and the edge
$x=a$ follows from the symmetric case $I_1\times I_0$. Thus all four
boundary edges of $I_0\times I_0$, namely $x=0$, $v=0$, $x=a$, and
$v=a$, have been treated. Since $G$ has no local maximum in the
interior, it follows that $G\leq d_1$ throughout
$I_0\times I_0$.






Finally, consider the feasible part of $I_0\times I_2$. Since
$x\in[0,a]$, $v\in[b,1]$, and $v\leq\phi(x)=h_K(x)$, it is described
by
\begin{equation*}
 0\leq x\leq a,
 \qquad
 b\leq v\leq h_K(x).
\end{equation*}
Recall that $h_K(0)=1$ and $h_K(a)=b$. Hence, the boundary of the
feasible region consists of
\begin{equation*}
 \{(0,v):b\leq v\leq1\},
 \qquad
 \{(x,b):0\leq x\leq a\},
 \qquad
 \{(x,h_K(x)):0\leq x\leq a\}.
\end{equation*}
When $x=a$, the inequalities
$
 b\leq v\leq h_K(a)=b
$
force $v=b$. Thus, $x=a$ contributes only the point $(a,b)$, which
already belongs to both the second and the third boundary parts.
Similarly, $v=1$ forces $x=0$, giving the point $(0,1)$, which already belongs to both the first and the third boundary parts.
The boundary $x=0$ was treated above, while the boundary $v=b$ was
treated in \eqref{eq:Qb-factor}. On the curved boundary
$v=h_K(x)=\phi(x)$, the involution property of $\phi$ gives
$
 \phi(v)=\phi(\phi(x))=x.
$
Therefore,
\begin{align*}
 G_{p,q}(x,v)
 =
 \bigl((1-p)x+p\phi(v)\bigr)
 \bigl((1-q)v+q\phi(x)\bigr)=xv=xh_K(x).
\end{align*}
Since $0\leq x\leq a\leq b$, we have
\begin{equation*}
 d_1-xh_K(x)
 =
 \frac{d_1(a-x)(b-x)}{d_1+cx}
 \geq0.
\end{equation*}
Thus, $G_{p,q}\leq d_1$ on every boundary part. The preceding
second-derivative argument shows that $G_{p,q}$ has no local maximum
in the interior. Hence, $G_{p,q}\leq d_1$ throughout the feasible part
of $I_0\times I_2$. As before, interchanging $(x,p)$ with $(v,q)$ maps the feasible part
of $I_2\times I_0$ to that of $I_0\times I_2$. Since
$
 G_{p,q}(x,v)=G_{q,p}(v,x),
$
the same argument applies to $I_2\times I_0$.
This completes the proof.
\end{proof}

\end{document}
